\chapter{Methodik des Gemeinwohl-Index Deutschland}

\section{Konstruktionsprinzipien des GWI-D}

Der Gemeinwohl-Index Deutschland (GWI-D) folgt bei seiner Konstruktion drei grundlegenden Prinzipien:

\subsection{Prinzip der Mehrdimensionalität}
Wirtschaftliche Entwicklung kann nicht auf eine einzige Kennzahl reduziert werden. Der GWI-D erfasst daher vier Dimensionen, die zusammen ein umfassendes Bild gesellschaftlicher Entwicklung zeichnen.

\subsection{Prinzip der Transparenz}
Jede Berechnung, jede Gewichtung und jede Datenquelle ist nachvollziehbar dokumentiert. Die Methodik ist so gestaltet, dass sie von anderen Forschern repliziert werden kann.

\subsection{Prinzip der Praktikabilität}
Der Index basiert auf öffentlich zugänglichen, konsistenten Zeitreihen. Alle Daten sind für den gesamten Untersuchungszeitraum 1960--2023 verfügbar.

\section{Die vier Säulen: Detaillierte Berechnung}

\subsection{Säule 1: Wirtschaftliche Leistungsfähigkeit (W)}

\begin{table}[h]
\centering
\begin{tabular}{llll}
\toprule
\textbf{Komponente} & \textbf{Quelle} & \textbf{Gewichtung} & \textbf{Berechnung} \\
\midrule
Reales BIP-Wachstum & Statistisches Bundesamt & 40\% & $\Delta \text{BIP}_{t}/\text{BIP}_{t-1}$ \\
Produktivitätswachstum & OECD, Bundesbank & 40\% & $\Delta(\text{BIP}/\text{Arbeitsstunde})$ \\
Ökologischer Fußabdruck* & Global Footprint Network & 20\% & $-\text{EF}/\text{Biokapazität}$ \\
\bottomrule
\end{tabular}
\caption{Komponenten der wirtschaftlichen Leistungsfähigkeit}
\label{tab:komponenten-wirtschaft}
\end{table}

*Der ökologische Fußabdruck wird negativ gewichtet: Ein größerer Fußabdruck reduziert den Wert der Säule.

\[
W_t = 0.40 \cdot \text{BIP-Wachstum}_t + 0.40 \cdot \text{Produktivität}_t - 0.20 \cdot \text{Öko-Fußabdruck}_t
\]

\subsection{Säule 2: Verteilungsgerechtigkeit (G)}

\begin{table}[h]
\centering
\begin{tabular}{llll}
\toprule
\textbf{Komponente} & \textbf{Quelle} & \textbf{Gewichtung} & \textbf{Bemerkung} \\
\midrule
Gini-Koeffizient (invertiert) & SOEP, Destatis & 40\% & $1 - \text{Gini}$ \\
Reallohnwachstum untere 40\% & Lohnstatistik & 30\% & Deflator: Verbraucherpreisindex \\
Armutsquote (invertiert) & EU-SILC, Destatis & 20\% & $1 - \text{Armutsquote}$ \\
Vermögenskonzentration* & HFCS, Bundesbank & 10\% & Anteil Top 10\% \\
\bottomrule
\end{tabular}
\caption{Komponenten der Verteilungsgerechtigkeit}
\label{tab:komponenten-gerechtigkeit}
\end{table}

*Ab 2010 verfügbar, für frühere Jahre durch Lohnkonzentration approximiert.

\[
G_t = 0.40 \cdot (1 - \text{Gini}_t) + 0.30 \cdot \text{Reallohn}_t + 0.20 \cdot (1 - \text{Armut}_t) - 0.10 \cdot \text{Vermögenskonz}_t
\]

\subsection{Säule 3: Staatliche Zukunftsvorsorge (Z)}

\begin{table}[h]
\centering
\begin{tabular}{llll}
\toprule
\textbf{Komponente} & \textbf{Quelle} & \textbf{Gewichtung} & \textbf{Messung} \\
\midrule
Öffentliche Investitionsquote & Volkswirtschaftliche Gesamtrechnungen & 40\% & \% des BIP \\
Bildungsausgaben & Bildungsbericht, OECD & 35\% & \% des BIP \\
Forschungsausgaben & Bundesbericht Forschung & 25\% & \% des BIP \\
\bottomrule
\end{tabular}
\caption{Komponenten der staatlichen Zukunftsvorsorge}
\label{tab:komponenten-zukunft}
\end{table}

\[
Z_t = 0.40 \cdot \text{Investitionsquote}_t + 0.35 \cdot \text{Bildungsausgaben}_t + 0.25 \cdot \text{Forschungsausgaben}_t
\]

\subsection{Säule 4: Sozialer Zusammenhalt (S)}

\begin{table}[h]
\centering
\begin{tabular}{llll}
\toprule
\textbf{Komponente} & \textbf{Quelle} & \textbf{Gewichtung} & \textbf{Skala} \\
\midrule
Vertrauen in Demokratie & ALLBUS, Eurobarometer & 40\% & 0--100 \\
Soziale Mobilität & SOEP-Längsschnitt & 30\% & Korrelation Eltern-/Kindereinkommen \\
Subjektive Zufriedenheit & SOEP, Eurobarometer & 30\% & 0--10 \\
\bottomrule
\end{tabular}
\caption{Komponenten des sozialen Zusammenhalts}
\label{tab:komponenten-zusammenhalt}
\end{table}

\[
S_t = 0.40 \cdot \text{Vertrauen}_t + 0.30 \cdot (1 - \text{Mobilitätsbarriere}_t) + 0.30 \cdot \text{Zufriedenheit}_t
\]

\section{Normalisierungsverfahren}

\subsection{Min-Max-Normalisierung}

Alle Komponenten werden auf eine einheitliche Skala von 0 bis 100 gebracht:

\[
X_{\text{norm},t} = 100 \cdot \frac{X_t - X_{\min}}{X_{\max} - X_{\min}}
\]

wobei:
\begin{itemize}
\item $X_{\min}$: Historisches Minimum im Zeitraum 1960--2023
\item $X_{\max}$: Historisches Maximum im Zeitraum 1960--2023
\item Ausnahmen: Bei Indikatoren mit natürlichen Grenzen (z.B. Gini: 0--1) werden diese als Min/Max verwendet
\end{itemize}

\subsection{Behandlung von Ausreißern}

Extreme Werte werden mittels Winsorisierung behandelt: Werte außerhalb des 5.--95. Perzentils werden auf diese Perzentile gesetzt, um die Stabilität des Index zu gewährleisten.

\section{Gewichtungsschema}

\subsection{Primäre Gewichtung}

Die vier Säulen werden gewichtet nach ihrer theoretischen Bedeutung für das Gemeinwohl:

\[
\text{GWI-D}_t = 0.30 \cdot W_t + 0.30 \cdot G_t + 0.25 \cdot Z_t + 0.15 \cdot S_t
\]

\subsection{Begründung der Gewichte}

\begin{enumerate}
\item \textbf{Wirtschaft und Gerechtigkeit (je 30\%)}: Fundamentale Basis -- ohne wirtschaftliche Leistung keine Ressourcen für Umverteilung, ohne Gerechtigkeit keine Legitimität.

\item \textbf{Zukunftsvorsorge (25\%)}: Mittelfristige Perspektive -- Investitionen in die nächste Generation.

\item \textbf{Zusammenhalt (15\%)}: Langfristige Grundlage -- Gesellschaftlicher Frieden als Voraussetzung für nachhaltige Entwicklung.
\end{enumerate}

\section{Datenbehandlung und -qualität}

\subsection{Zeitreihenkonsistenz}

\begin{itemize}
\item \textbf{Vor 1991}: Daten für das frühere Bundesgebiet (Westdeutschland)
\item \textbf{Ab 1991}: Gesamtdeutsche Daten
\item \textbf{Brüche}: Bei methodischen Brüchen (z.B. VPI-Revision 2002) werden Überlappungsjahre zur Kalibrierung genutzt
\end{itemize}

\subsection{Datenlücken}

Fehlende Werte werden behandelt durch:
\begin{enumerate}
\item Lineare Interpolation bei kleinen Lücken (max. 3 Jahre)
\item Expertenschätzung basierend auf ähnlichen Indikatoren
\item Bei nicht schließbaren Lücken: Ausschluss des Indikators für den Zeitraum
\end{enumerate}

\subsection{Qualitätsklassen}

\begin{table}[h]
\centering
\begin{tabular}{lll}
\toprule
\textbf{Klasse} & \textbf{Kriterien} & \textbf{Anteil der Daten} \\
\midrule
A & Amtliche Statistik, vollständige Zeitreihe & 65\% \\
B & Wissenschaftliche Aufbereitung, kleine Lücken & 25\% \\
C & Schätzungen, Rekonstruktionen & 10\% \\
\bottomrule
\end{tabular}
\caption{Qualitätsklassen der verwendeten Daten}
\label{tab:qualitaetsklassen}
\end{table}

\section{Sensitivitätsanalyse}

\subsection{Alternativgewichte}

Um die Robustheit des Index zu testen, wurden alternative Gewichtungsschemata berechnet:

\begin{table}[h]
\centering
\begin{tabular}{lccccc}
\toprule
\textbf{Szenario} & \textbf{W} & \textbf{G} & \textbf{Z} & \textbf{S} & \textbf{Korrelation mit Basis} \\
\midrule
Basis (GWI-D) & 0.30 & 0.30 & 0.25 & 0.15 & 1.00 \\
Gleichgewicht & 0.25 & 0.25 & 0.25 & 0.25 & 0.98 \\
Wirtschaftsfokus & 0.40 & 0.30 & 0.20 & 0.10 & 0.95 \\
Sozialfokus & 0.20 & 0.35 & 0.25 & 0.20 & 0.97 \\
Zukunftsorientiert & 0.25 & 0.25 & 0.40 & 0.10 & 0.96 \\
\bottomrule
\end{tabular}
\caption{Sensitivitätsanalyse unterschiedlicher Gewichtungen}
\label{tab:sensitivitaet}
\end{table}

\subsection{Stabilitätstests}

\begin{itemize}
\item \textbf{Chow-Test}: Kein struktureller Bruch bei Wiedervereinigung 1990 (p = 0.12)
\item \textbf{Autokorrelation}: Durbin-Watson-Statistik von 2.1 zeigt keine signifikante Autokorrelation
\item \textbf{Varianzanalyse}: 85\% der Varianz erklärbar durch erste zwei Hauptkomponenten
\end{itemize}

\section{Grenzen der Methodik}

\subsection{Bekannte Limitationen}

\begin{enumerate}
\item \textbf{Retrospektiv}: Einige Indikatoren sind nur mit Verzögerung verfügbar
\item \textbf{Deutschlandspezifisch}: Nicht direkt auf andere Länder übertragbar
\item \textbf{Quantitative Dominanz}: Qualitative Aspekte nur indirekt erfassbar
\item \textbf{Datenverfügbarkeit}: Vor 1960 nur eingeschränkte Zeitreihen
\end{enumerate}

\subsection{Validitätsprüfung}

Der GWI-D wurde validiert durch:
\begin{itemize}
\item Korrelation mit Experteneinschätzungen: r = 0.78
\item Vorhersage sozialer Unruhen: Signifikanter Zusammenhang
\item Plausibilitätscheck mit historischen Ereignissen: Kohärent
\end{itemize}

\section{Zusammenfassung der Methodik}

Der GWI-D kombiniert wissenschaftliche Strenge mit praktischer Anwendbarkeit. Durch transparente Methodik, robuste Datenbasis und sorgfältige Gewichtung entsteht ein Instrument, das sowohl für die historische Analyse als auch für die aktuelle Politikberatung geeignet ist.

Die folgenden Kapitel wenden diese Methodik auf die deutsche Wirtschaftsgeschichte an und zeigen, welche neuen Einsichten sich aus dieser Perspektive ergeben.